\documentclass[11pt,letterpaper]{article}
\usepackage{math-template}

\newcommand{\StudentName}{Your Name}
\newcommand{\CourseName}{MATH 101 --- Introduction to Analysis}
\newcommand{\AssignmentTitle}{Lecture Notes: Sequences}
\newcommand{\DueDate}{October 15, 2026}
\newcommand{\StudentNumber}{}
\newcommand{\TutorialSection}{}
\renewcommand{\DateLabel}{Date}

\begin{document}
\makeassignmentheader

\section{Convergence}
\begin{definition}[Convergence]
A real sequence $(a_n)$ converges to $L\in\R$ if for every $\eps>0$
there exists $N\in\N$ such that $n\geq N$ implies $\abs{a_n-L}<\eps$.
\end{definition}

\begin{example}
The sequence $a_n=1/n$ converges to zero. Given $\eps>0$, choose
$N>1/\eps$; then $n\geq N$ implies $1/n<\eps$.
\end{example}

\begin{theorem}[Uniqueness of limits]\label{thm:unique}
A convergent real sequence has exactly one limit.
\end{theorem}
\begin{proof}
Suppose $a_n\to L$ and $a_n\to M$ with $L\neq M$.
Set $\eps=\abs{L-M}/3$. For sufficiently large $n$,
\[
  \abs{L-M}\leq\abs{L-a_n}+\abs{a_n-M}
  <2\eps=\frac23\abs{L-M},
\]
a contradiction. Thus $L=M$.
\end{proof}

\begin{remark}
Theorem~\ref{thm:unique} justifies writing ``the limit'' of a convergent sequence.
\end{remark}

\section{A useful consequence}
\begin{proposition}
Every convergent real sequence is bounded.
\end{proposition}
\begin{proof}
Suppose $a_n\to L$. Choose $N$ so that $n\geq N$ implies
$\abs{a_n-L}<1$, and set
\[
  B=\max\set{\abs{a_1},\ldots,\abs{a_N},\abs{L}+1}.
\]
For $n<N$ the bound $\abs{a_n}\leq B$ follows from the definition of $B$.
For $n\geq N$, the triangle inequality gives
$\abs{a_n}\leq\abs{a_n-L}+\abs{L}<\abs{L}+1\leq B$.
\end{proof}
\end{document}
